\section{Bayesian 线性回归把系数不确定性传到预测}
\label{sec:13-Machine-Learning/02-Linear-Models/04}

\subsection{评审问，这条误差棒管的是谁}

一个门店备货模型上线前要过评审。拟合的是 \(y=x^{\trans}\beta+\varepsilon\)，最小二乘给出 \(\widehat\beta\)，新门店的特征 \(x_*\) 代进去得到点预测 458 件，报告上配了一条 \(\pm 2\SE\) 的区间。评审问这条棒子管的是谁：是这家新店下个月的实际销量，还是别的什么。

答案让人不太舒服。\(\SE\bigl(x_*^{\trans}\widehat\beta\bigr)\) 描述的是把整个采数据加拟合的流程重复无数次时，那条直线在 \(x_*\) 处的高度会怎样摆动；它既不包含这家店自己的随机波动，也不是对``这一笔预测''的信念陈述。备货的人想要的是后者，而经典区间给的是前者（两种区间各自在说什么，见\hyperref[ch:089]{``区间估计：让不确定性进入答案''}）。

要把``我对这家店的销量相信到什么程度''直接写出来，得先给系数一个概率分布，再在预测时把系数积掉。

\subsection{两份精度直接相加}

先假定噪声方差 \(\sigma^2\) 已知，把一层不确定性搁到后面。模型是
\[
y\given\beta\sim\mathcal N(X\beta,\sigma^2I),
\qquad
\beta\sim\mathcal N(m_0,S_0),\quad S_0\succ 0,
\]
设计矩阵 \(X\) 视为给定。\(m_0\) 是看数据之前对系数最合理的猜测，\(S_0\) 表达对这个猜测的把握。

后验的骨架是两个二次型相加再配方。丢掉与 \(\beta\) 无关的常数，
\[
\log p(\beta\given y,X)=-\frac{1}{2\sigma^2}\norm{y-X\beta}_2^2-\frac12(\beta-m_0)^{\trans}S_0^{-1}(\beta-m_0)+C .
\]
两项都是 \(\beta\) 的二次函数，和仍是二次函数，指数上是二次函数的密度只能是 Gaussian。展开收集二次项与一次项，得到
\[
\Lambda_n=S_0^{-1}+\sigma^{-2}X^{\trans}X,\qquad
h_n=S_0^{-1}m_0+\sigma^{-2}X^{\trans}y,\qquad
\beta\given y,X\sim\mathcal N\bigl(\Lambda_n^{-1}h_n,\ \Lambda_n^{-1}\bigr).
\]
条件对账：用到 \(S_0\) 正定（否则 \(S_0^{-1}\) 不存在，先验也不是 proper 的），用到噪声协方差是 \(\sigma^2I\) 且 \(\sigma^2\) 已知，用到 \(X\) 不随 \(\beta\) 变化。没有用到 \(X^{\trans}X\) 可逆——这一点下面还要单独讲。共轭性为什么让更新保持在同一族里，见\hyperref[sec:13-Machine-Learning/06-Probabilistic-Models/07]{``共轭先验让 Bayesian 更新保持同一形式''}。

精度相加这条式子值得盯一会儿。先验精度和数据精度按方向叠加，数据只在 \(X^{\trans}X\) 有能量的方向上提高精度；样本没有覆盖到的方向，后验精度仍旧等于先验精度。先验不是可以随手设掉的初始条件，它决定了数据沉默的方向上后验停在哪里。

\subsection{Ridge 只用了后验的一个点}

后验众数是最小化
\[
\frac{1}{2\sigma^2}\norm{y-X\beta}_2^2+\frac12(\beta-m_0)^{\trans}S_0^{-1}(\beta-m_0)
\]
的那个 \(\beta\)。取 \(m_0=0\)、\(S_0=\tau^2I\)，同乘 \(2\sigma^2\) 之后正是 \(\norm{y-X\beta}_2^2+\lambda\norm{\beta}_2^2\)，其中 \(\lambda=\sigma^2/\tau^2\)。上一节的 Ridge 就是这个先验下的 MAP。

这个对应有边界。给不同系数不同的先验方差、或者对截距放宽先验，MAP 就不再是标准 Ridge；上一节提到的``换了单位同一个 \(\lambda\) 含义就变''，在这里读作先验尺度必须跟着特征尺度走。更要紧的是 MAP 只是一个点，它把 \(S_n=\Lambda_n^{-1}\) 整个丢掉了。协方差矩阵记录的是系数之间的联合不确定：哪些线性组合被数据钉死了，哪些组合仍然可以自由滑动。共线时，\(S_n\) 在 \(\beta_1-\beta_2\) 方向上有一个大的特征值，这正是上一节那个``分不清谁贡献了多少''的定量表述。

\subsection{预测方差分成两段，只有一段会随样本消失}

对新特征 \(x_*\)，模型规定 \(Y_*=x_*^{\trans}\beta+\varepsilon_*\)。把 \(\beta\) 按后验积掉——两个 Gaussian 的线性组合仍是 Gaussian——得到
\[
Y_*\given x_*,X,y\sim\mathcal N\bigl(x_*^{\trans}m_n,\ \sigma^2+x_*^{\trans}S_nx_*\bigr).
\]
方差的两段各有出处：\(\sigma^2\) 是这家店自己的随机波动，\(x_*^{\trans}S_nx_*\) 是我们对 \(\beta\) 的无知在 \(x_*\) 方向上的投影。同分布的数据越多，后者越小；前者不动。样本量能买到对系数的了解，买不到对世界的确定。

\begin{center}
\begin{tikzpicture}[>=stealth]
  \fill[black!12]
    plot[domain=-2.4:2.4,samples=60] ({\x},{0.32*\x+sqrt(0.16+0.10*\x*\x)})
    -- plot[domain=2.4:-2.4,samples=60] ({\x},{0.32*\x-sqrt(0.16+0.10*\x*\x)})
    -- cycle;
  \draw[->,black!60] (-2.8,-1.55) -- (3.0,-1.55) node[below] {\footnotesize \(x\)};
  \draw[dashed,black!70] (-2.4,-1.168) -- (2.4,1.168);
  \draw[dashed,black!70] (-2.4,-0.368) -- (2.4,1.968);
  \draw[very thick] (-2.4,-0.768) -- (2.4,1.568);
  \fill (-1.2,-0.31) circle (1.5pt);
  \fill (-0.8,-0.34) circle (1.5pt);
  \fill (-0.35,0.02) circle (1.5pt);
  \fill (0.15,-0.02) circle (1.5pt);
  \fill (0.6,0.34) circle (1.5pt);
  \fill (1.05,0.22) circle (1.5pt);
  \draw[black!60] (-1.3,-1.42) -- (-1.3,-1.55) -- (1.3,-1.55) -- (1.3,-1.42);
  \node[black!60,below] at (0,-1.62) {\footnotesize 训练数据覆盖范围};
  \node[right] at (2.45,1.97) {\footnotesize \(\pm 2\sigma\)};
  \node[right] at (2.45,2.55) {\footnotesize 含参数不确定};
  \draw[black!45] (2.42,2.5) -- (2.15,2.15);
\end{tikzpicture}
\end{center}

\emph{虚线带只含观测噪声，宽度处处一样；阴影带还含系数不确定，离开数据覆盖的方向就张开。}

图里两条虚线相距 \(4\sigma\)，与 \(x\) 无关。阴影带在数据密集处几乎贴着虚线，往两端逐渐撑开，撑开的那部分全部来自 \(x_*^{\trans}S_nx_*\)。这条形状是 Bayesian 线性回归相对点估计最实在的收益：模型在自己没见过的区域会主动把话说软。

一组新样本 \(X_*\) 的联合预测是 \(\mathcal N\bigl(X_*m_n,\ \sigma^2I+X_*S_nX_*^{\trans}\bigr)\)，非对角元素不为零。逐点报告边际区间会丢掉这个结构：若后验里 \(\beta\) 偏高，一整批门店的预测会同时偏高，而各自的边际区间看不出这种同向移动。总量的区间要从联合分布里算，不能把逐点区间简单相加。

\subsection{远近由 \(S_n\) 定义，不由欧氏距离定义}

\(x_*^{\trans}S_nx_*\) 是一个二次型，它衡量的不是 \(x_*\) 离训练数据中心有多远，而是 \(x_*\) 指向哪些方向。落在训练设计充分激活的方向上，\(S_n\) 在那里已经被压得很小，区间不会怎么变宽；落在设计几乎没有覆盖的方向上——比如一个从未同时出现过的特征组合——同样的欧氏距离会换来大得多的方差。``离数据远''这句话应当在 \(S_n\) 诱导的度量下理解。

这也意味着先验尺度必须与特征尺度一起被记录和版本管理。特征标准化之后，各向同性先验对所有方向施加同等收缩；不标准化时，同一个 \(\tau\) 对以元为单位和以万元为单位的两个变量施加的实际约束差了四个数量级。换单位不该悄悄改变推断结论。

外推时还有一层区间盖不住的东西。就算后验预测带在边缘张得很开，它仍然建立在``均值确实是线性的''这个前提上。真实关系若有弯曲，再宽的带也可能整段偏在一侧。Bayesian 给出的是给定模型之后的条件不确定性，模型形式本身错了的那部分无知，不在这个积分里。

\subsection{先验补上秩亏方向，但那不叫学到了}

\(p>n\) 或特征线性相关时 \(X^{\trans}X\) 秩亏，最小二乘的解不唯一。加了 proper 先验之后 \(\Lambda_n=S_0^{-1}+\sigma^{-2}X^{\trans}X\) 一定可逆，后验唯一，预测分布照样算得出来。这是数值上的好消息，读解释时却要克制。

在那些秩亏方向上，后验精度全部来自 \(S_0^{-1}\)，也就是说后验等于先验，数据一个字也没说。报告后验均值时若不区分哪些方向被数据推动过、哪些方向只是先验的回声，读者会把先验当成结论。一个可操作的检查是比较先验方差与后验方差沿各个方向的比值：比值接近 1 的方向，数据没有贡献。这与上一节判断共线的做法是同一件事换了语言。

\subsection{\(\sigma^2\) 未知时，尾巴自己变厚}

现实里 \(\sigma^2\) 当然未知。给 \(\sigma^2\) 一个 inverse-gamma 先验、并让 \(\beta\) 在给定 \(\sigma^2\) 时为 Gaussian，这个 normal--inverse-gamma 族仍然共轭，系数与噪声水平可以一起更新。把两者都积掉之后，后验预测不再是 Gaussian 而是多元 Student-t，自由度随样本量增长。

厚尾不是修辞。把 \(\sigma^2\) 固定成残差 MSE 这个点估计，相当于假装噪声水平已经知道，有限样本下这会低估极端情形的概率；Student-t 的尾部把``连噪声多大都还没定''这层不确定折算了进来。若给 \(\sigma^2\) 用很弱的 improper 先验，得先验证联合后验是不是 proper——样本不足或设计秩亏时积分可能发散，无信息先验不等于没有假设。

\subsection{实现里不出现逆矩阵}

\(S_n=\Lambda_n^{-1}\) 是写给人看的。算的时候对 \(\Lambda_n\) 做 Cholesky 分解 \(\Lambda_n=LL^{\trans}\)，后验均值由 \(\Lambda_nm_n=h_n\) 回代求得，采样用 \(m_n+L^{-\trans}z\)，对数行列式由 \(2\sum_i\log L_{ii}\) 得到，全程不形成逆矩阵（理由见\hyperref[sec:08-Numerical-Analysis/07-Numerical-Methods-for-ML/03]{``需要的是解，而不是逆矩阵''}与\hyperref[sec:08-Numerical-Analysis/05-Numerical-Linear-Algebra/03]{``Cholesky：利用正定结构''}）。\(p\gg n\) 时用 Woodbury 恒等式
\[
\bigl(\sigma^{-2}X^{\trans}X+S_0^{-1}\bigr)^{-1}
=S_0-S_0X^{\trans}\bigl(\sigma^2I+XS_0X^{\trans}\bigr)^{-1}XS_0
\]
把主要运算搬到 \(n\times n\) 空间。等号右边的形式已经只通过 \(XS_0X^{\trans}\) 依赖数据，把它换成一个核矩阵就得到 Gaussian process 回归（见\hyperref[sec:13-Machine-Learning/09-Kernel-Methods-and-Gaussian-Processes/05]{``Gaussian process 把核变成函数先验''}）。

最后一步是验区间而不是只画区间。后验预测检查的做法是从后验预测里生成若干套假数据，看它们的残差分布、极值和分组均值与真实数据是否像；再在独立数据上数一遍名义 \(95\%\) 区间的实际覆盖率。覆盖率系统性偏低，说明似然或先验哪里错了，加迭代次数和提高浮点精度都修不好。

本节的先验、似然和链接方式都锁定在 Gaussian 上，换一种响应类型就得整套重来。下一节把这件事做成通则：响应分布、连接函数和线性预测子各司其职，前面几节的线性回归与 logistic 回归会作为同一张表里的两行出现。
